
% ===== DRONES TERMS =====
\dictentry{2-Axis Gimbal}{A camera gimbal that stabilizes footage along two rotational axes: pitch (tilt) and roll (bank). By counteracting the drone's forward/backward tilt and side-to-side banking, it keeps the camera level during forward flight. It is lighter and cheaper than a 3-axis gimbal but cannot compensate for yaw rotation, making it suitable for fixed-wing platforms.}{catD}
\dictentry{2.4 GHz}{The 2.4 GHz ISM band is the most widely used frequency for drone radio control and telemetry links. It offers a good balance between range, data throughput, and antenna size, making it ideal for compact receivers. Most modern RC protocols such as CRSF and ELRS operate on this band with frequency-hopping to avoid interference.}{catD}
\dictentry{3-Axis Gimbal}{A camera gimbal stabilizing pitch, roll, and yaw for fully stabilized footage. By correcting rotation on all three axes, it ensures a smooth, level video feed even during aggressive maneuvers. This is the standard gimbal type on most modern camera drones for professional aerial cinematography.}{catD}
\dictentry{433 MHz}{The 433 MHz band is an ultra-long-range frequency used for drone telemetry and control links. Its lower frequency provides superior signal penetration and longer range compared to 2.4 GHz, but with lower data throughput. It is commonly used for long-range flight missions where link reliability at distance is critical.}{catD}
\dictentry{5.8 GHz}{The 5.8 GHz ISM band is the primary frequency for analog and digital FPV video transmission. It provides higher bandwidth than 2.4 GHz, allowing for clearer video with less latency, though at shorter effective range. Most FPV transmitters and goggles operate on this band with multiple channels to avoid mutual interference.}{catD}
\dictentry{900 MHz}{The 900 MHz band is a longer-range frequency used for drone telemetry and control links. Its lower frequency provides better signal propagation and obstacle penetration compared to higher bands. It is commonly used for telemetry modules and long-range RC protocols where maintaining link at great distances is essential.}{catD}
\dictentry{Accel Calibration}{The process of calibrating the flight controller's accelerometer to ensure accurate acceleration readings. During calibration, the drone must be placed on a perfectly level surface so the sensor can establish its zero-g reference point. Incorrect calibration leads to drift, unstable hovering, and unreliable altitude hold in stabilized flight modes.}{catD}
\dictentry{Accelerometer}{A micro-electromechanical sensor that measures linear acceleration along three orthogonal axes. It detects gravitational pull and translational movement, providing essential data for the flight controller's stabilization algorithms. Combined with the gyroscope, it forms the core of the inertial measurement unit (IMU) that keeps drones level and responsive to pilot input.}{catD}
\dictentry{Aerial Photography}{The practice of capturing images or video from a camera mounted on a drone. Drones provide unique elevated perspectives that were previously only achievable with manned aircraft or helicopters. Modern drones use stabilized gimbals and high-resolution cameras to produce professional-quality aerial images for real estate, events, and journalism.}{catD}
\dictentry{Agriculture Drone}{A drone optimized for crop monitoring, mapping, and precision spraying of pesticides or fertilizers. These drones carry multispectral cameras or liquid tanks and fly pre-programmed patterns over fields to treat crops efficiently. They reduce chemical usage and labor costs while enabling farmers to identify problem areas through aerial imagery analysis.}{catD}
\dictentry{AirSim}{A drone and car simulator developed by Microsoft using Unreal Engine for photorealistic rendering. It provides a high-fidelity physics engine and realistic sensor simulation for training autonomous vehicles without real-world risk. AirSim supports both Windows and Linux and interfaces with ROS, making it a popular tool for research and drone AI development.}{catD}
\dictentry{Airspace}{The three-dimensional volume of atmosphere above the earth's surface in which drones operate, classified into controlled and uncontrolled zones. Different classes of airspace carry different requirements for altitude, equipment, and pilot certification. Drone pilots must understand airspace classification to determine where and when they can legally fly.}{catD}
\dictentry{Altitude Hold}{A flight mode in which the drone maintains a constant barometric altitude without pilot throttle input. The barometer provides relative altitude data to the flight controller, which adjusts motor output to counteract vertical drift. This mode is essential for stable aerial photography and surveying operations where consistent height is required.}{catD}
\dictentry{Angle Mode}{A beginner-friendly flight mode that limits the maximum tilt angle of the drone for stable, predictable flight. When the pilot releases the sticks, the drone automatically levels itself back to horizontal. It provides a safety net for new pilots but limits agility, making it unsuitable for aerobatic or racing maneuvers.}{catD}
\dictentry{Antenna}{A device that transmits or receives radio signals for drone communication links, including control, telemetry, and video. Antenna design determines the range, directionality, and reliability of the radio link between the pilot and the aircraft. Different antenna types such as omnidirectional, directional, and patch are chosen based on the mission requirements.}{catD}
\dictentry{Antenna Polarization}{The orientation of the electromagnetic wave emitted by an antenna, which must match between the transmitter and receiver for optimal signal strength. Common types include linear (vertical or horizontal) and circular (left-hand or right-hand) polarization. Mismatched polarization between the transmitting and receiving antennas can cause significant signal loss of 20 dB or more.}{catD}
\dictentry{Anti-Vibration Mount}{A mounting system that dampens vibrations generated by the drone's motors and propellers before they reach sensitive sensors. It typically uses rubber grommets, foam pads, or silicone dampeners to isolate the flight controller and IMU from high-frequency oscillations. Proper vibration isolation is critical for accurate sensor readings and stable flight performance.}{catD}
\dictentry{ArduPilot}{An open-source autopilot software suite supporting a wide range of unmanned vehicles including multirotors, fixed-wing aircraft, and rovers. It runs on various flight controller hardware and provides features like autonomous missions, geofencing, and extensive sensor integration. ArduPilot has one of the largest communities in the open-source drone ecosystem with thousands of contributors worldwide.}{catD}
\dictentry{ArduPilot Project}{The global community of developers, testers, and users who collaborate on developing and maintaining the ArduPilot open-source autopilot software. The project operates through GitHub with contributions from hobbyists, researchers, and commercial drone companies alike. It provides extensive documentation, forums, and flight test resources that support drone builders of all skill levels.}{catD}
\dictentry{Arm}{A structural extension of the drone's frame that carries a motor and propeller at its end. Arms connect the central body to the propulsion units and determine the drone's overall size and motor configuration. Their length and material affect the drone's flight characteristics, with longer arms providing greater leverage but adding weight.}{catD}
\dictentry{Arming}{The process of enabling motor power on the drone, typically requiring a specific stick combination or switch on the transmitter as a safety precaution. Arming activates the electronic speed controllers and allows the motors to spin, so the drone is ready for takeoff. Disarming immediately stops all motors and is used as a safety measure during ground handling or after landing.}{catD}
\dictentry{ASTM International}{An international standards organization that develops and publishes technical standards for drone design, safety, and operational procedures. ASTM standards cover areas such as unmanned aircraft systems performance, counter-UAS technologies, and delivery drone operations. Compliance with ASTM standards is often required or referenced by regulatory agencies for commercial drone certification.}{catD}
\dictentry{Autonomous Drone}{A drone capable of navigating and executing missions independently without continuous pilot input. It uses onboard sensors, GPS, and artificial intelligence to make real-time decisions about flight path, obstacle avoidance, and mission objectives. Autonomous drones are used in applications ranging from agricultural surveying to delivery services where human piloting is impractical.}{catD}
\dictentry{Autotune}{An automatic PID tuning procedure performed by the flight controller to optimize its response characteristics for a specific drone configuration. During autotune, the controller applies small control inputs and measures the drone's response to calculate optimal proportional, integral, and derivative gains. This process saves significant time compared to manual tuning and typically produces good baseline settings for most flight scenarios.}{catD}
\dictentry{Barometer}{A sensor that measures atmospheric pressure to estimate the drone's altitude relative to a reference point. It provides essential data for altitude hold mode and flight controller stabilization, but is sensitive to temperature changes and wind effects. Barometric altitude is relative rather than absolute, so it requires calibration at takeoff for accurate readings throughout the flight.}{catD}
\dictentry{Battery Capacity}{The total amount of energy stored in a battery, measured in milliampere-hours (mAh) for drone applications. Higher capacity batteries provide longer flight times but add weight, requiring a careful balance between endurance and performance. A 5000 mAh battery, for example, can theoretically supply 5000 milliamps for one hour, though actual flight time depends on the drone's power consumption.}{catD}
\dictentry{Battery Failsafe}{An automatic action triggered by the flight controller when the battery voltage or capacity drops below a predefined safety threshold. The failsafe typically initiates a return-to-home sequence, a controlled landing, or a warning to the pilot depending on the severity of the voltage drop. This feature prevents catastrophic power loss during flight, which could lead to an uncontrolled crash.}{catD}
\dictentry{Beacon}{An audible or visual signaling device mounted on a drone to assist in locating it after a crash or emergency landing. Beacons may emit flashing LEDs, sounds, or RF signals that can be detected by the pilot or ground equipment. They are especially useful for finding drones that have gone down in tall grass, dense foliage, or low-visibility conditions.}{catD}
\dictentry{BEC}{A Battery Eliminator Circuit that steps down the main battery voltage to a regulated lower voltage for powering the flight controller, receiver, and other electronics. By providing a clean and stable power supply, the BEC eliminates the need for a separate receiver battery, reducing weight and complexity. Many ESCs include a built-in BEC, though standalone units are used for higher power requirements.}{catD}
\dictentry{Betaflight}{A popular open-source firmware optimized for racing and freestyle FPV drones, known for its highly responsive flight characteristics. It supports a wide range of flight controllers and provides extensive tuning options for PID controllers, filters, and flight modes. Betaflight has a large community of developers and users who continuously improve performance and add new features.}{catD}
\dictentry{Betaflight OSD}{The On-Screen Display system integrated into Betaflight firmware that overlays flight data onto the FPV video feed. It shows critical information such as battery voltage, flight time, signal strength, and artificial horizon directly in the pilot's goggles or monitor. The OSD configuration is customizable through the Betaflight configurator, allowing pilots to choose which data fields to display.}{catD}
\dictentry{Bidirectional DShot}{An advanced digital ESC protocol that allows two-way communication between the flight controller and the electronic speed controller. It not only sends throttle commands to the ESC but also receives telemetry data such as motor RPM, temperature, and current draw back to the flight controller. This feedback enables features like RPM-based dynamic filtering and more precise motor control for smoother flight.}{catD}
\dictentry{Bitrate}{The amount of data transmitted per second in a video stream, typically measured in megabits per second (Mbps). Higher bitrates produce clearer video with fewer compression artifacts but require more bandwidth and increase transmission latency. FPV pilots must balance bitrate against link reliability, since excessively high bitrates may cause signal breakup at longer ranges.}{catD}
\dictentry{Blackbox}{A flight data logging system that records sensor readings, motor outputs, pilot inputs, and other internal variables during flight. The recorded data can be analyzed using specialized tools to diagnose flight issues, tune PID settings, and troubleshoot vibrations. Blackbox logging is essential for advanced tuning and is supported by most modern flight controller firmware.}{catD}
\dictentry{Brushless Motor}{An electric motor that uses electronic commutation instead of physical brushes, resulting in higher efficiency, greater reliability, and longer lifespan. Brushless motors provide superior power-to-weight ratios and are the standard propulsion choice for modern drones of all sizes. They require an electronic speed controller (ESC) to manage the timing and sequence of electrical pulses to the motor windings.}{catD}
\dictentry{ButterFlight}{A fork of Betaflight firmware designed for exceptionally smooth and locked-in flight characteristics. It emphasizes reduced noise, refined filtering, and optimized PID algorithms to deliver a creamy flight feel favored by freestyle pilots. Though development has largely been absorbed back into Betaflight, ButterFlight's innovations influenced filtering and tuning approaches across the FPV community.}{catD}
\dictentry{Buzzer}{A small audible device mounted on the drone that emits sound to indicate status changes or assist in locating the aircraft. Buzzers are used for low-battery warnings, arming status confirmation, and lost-model alarms when the drone needs to be found after a crash. They are typically activated through a switch on the transmitter or automatically by the flight controller's failsafe routines.}{catD}
\dictentry{BVLOS}{Beyond Visual Line of Sight, an operational mode where the drone flies outside the pilot's direct visual observation range. BVLOS operations require special regulatory waivers and additional safety equipment such as detect-and-avoid systems and redundant communication links. This capability is critical for long-range missions including pipeline inspection, wildlife surveys, and delivery operations.}{catD}
\dictentry{C Rating}{The discharge rate rating of a LiPo battery that indicates how quickly it can safely deliver its stored energy relative to its capacity. A 100C-rated 1500 mAh battery can theoretically deliver 150 amps continuously without damage. Higher C ratings allow for more aggressive throttle bursts without voltage sag, which is especially important for racing and freestyle drones.}{catD}
\dictentry{Calibration}{The process of setting sensor offsets and scales on the flight controller to ensure accurate and reliable data readings. Calibration includes accelerometer zeroing, gyroscope bias correction, compass hard and soft iron adjustments, and radio channel endpoint configuration. Proper calibration is essential for stable flight and accurate autonomous navigation.}{catD}
\dictentry{Camera Mount}{The bracket or attachment system that secures the camera to the drone's frame at the desired angle and position. Camera mounts range from simple fixed-angle brackets to vibration-dampened quick-release systems for different shooting scenarios. The mount angle determines the camera's field of view during forward flight, with steeper angles preferred for faster drones.}{catD}
\dictentry{Camera Payload}{A camera system mounted on a drone specifically for aerial imaging, mapping, or surveillance missions. Camera payloads range from simple action cameras to professional-grade systems with zoom, thermal, or multispectral capabilities. The choice of camera payload depends on the mission requirements, including resolution, field of view, and spectral range needed.}{catD}
\dictentry{CAN Bus}{Controller Area Network, a robust serial communication bus used to connect sensors, ESCs, and other peripherals to the flight controller. CAN Bus provides reliable, high-speed data transfer with error detection, making it ideal for critical drone subsystems. It reduces wiring complexity compared to individual signal cables and supports hot-swapping of compatible devices.}{catD}
\dictentry{Cargo Pod}{An enclosed container mounted on a delivery drone for transporting packages, medical supplies, or other goods. Cargo pods are designed to securely hold payloads during flight while providing a mechanism for controlled release at the destination. They must account for weight distribution, aerodynamic drag, and quick-release systems to ensure safe and efficient delivery operations.}{catD}
\dictentry{Cinematic Drone}{A drone specifically designed or configured for high-quality aerial filmmaking, featuring smooth flight characteristics and professional camera systems. Cinematic drones prioritize stable, slow movements and long flight times over speed or agility, often carrying heavy camera payloads. They are used in film production, television, and commercial advertising to capture sweeping aerial shots.}{catD}
\dictentry{Circular Polarization}{A type of antenna polarization in which the electromagnetic wave rotates as it propagates through space. Circular polarization is more tolerant of antenna orientation changes than linear polarization, making it ideal for FPV video where the drone's orientation constantly changes. It comes in two variants, left-hand (LHCP) and right-hand (RHCP), which must match between transmitter and receiver.}{catD}
\dictentry{Cleanflight}{An early open-source drone flight controller firmware that was one of the first to bring advanced stabilization to hobbyist drones. It served as the foundation for later firmwares like Betaflight and INAV, introducing concepts like configurable PID controllers and modular sensor support. While largely superseded, Cleanflight's architecture influenced the design of modern FPV firmware.}{catD}
\dictentry{Codec}{A video compression algorithm that reduces raw video data size for efficient wireless transmission. Common drone codecs include H.264 and H.265, which achieve high compression ratios while maintaining acceptable image quality. The choice of codec directly impacts video latency, image clarity, and bandwidth requirements for the FPV video link.}{catD}
\dictentry{Collision Avoidance}{An onboard system that detects obstacles in the drone's flight path and takes evasive action to prevent crashes. It uses sensors such as stereo cameras, ultrasonic rangefinders, or time-of-flight sensors to perceive the environment. Modern systems can detect obstacles at ranges of 10 to 30 meters and automatically brake, hover, or reroute around the obstruction.}{catD}
\dictentry{Collision Avoidance Swarm}{A system that prevents collisions between individual drones operating in close proximity within a swarm formation. It combines inter-drone communication, relative position tracking, and distributed path planning to maintain safe separation distances. Each drone dynamically adjusts its flight path based on the positions and velocities of its neighbors to prevent mid-air collisions.}{catD}
\dictentry{Combat Drone}{An armed unmanned aerial vehicle designed for military strike, close air support, or kamikaze attack missions. Combat drones range from small loitering munitions to large aircraft carrying precision-guided ordnance. They reduce risk to human pilots while providing persistent surveillance and rapid response capabilities on the modern battlefield.}{catD}
\dictentry{Compass}{A magnetic heading sensor that detects the Earth's magnetic field to provide directional reference for the flight controller. It enables the drone to know its heading relative to magnetic north, which is essential for waypoint navigation, return-to-home functions, and position hold. Compass accuracy is affected by nearby magnetic interference from motors, batteries, and metal structures.}{catD}
\dictentry{Compass Calibration}{The process of calibrating the drone's magnetometer to compensate for hard iron and soft iron magnetic distortions from the drone's own structure. During calibration, the drone is rotated through all orientations to map the local magnetic field and calculate correction offsets. Failure to calibrate the compass properly can cause erratic heading, flyaways, and unreliable GPS-assisted flight modes.}{catD}
\dictentry{Compass Error}{A failsafe condition triggered when the flight controller detects unreliable or conflicting compass data during flight. The error may result from electromagnetic interference, sensor failure, or poor calibration quality. When a compass error is detected, the flight controller typically disables GPS-dependent modes like position hold and return to home as a safety precaution.}{catD}
\dictentry{Compass Mount}{A mounting system that isolates the drone's compass from electromagnetic interference generated by motors, ESCs, and power wiring. Compass mounts often use a raised mast or vibration-dampened platform to distance the sensor from the main electronics. Proper compass mounting is essential for accurate heading data, as even small amounts of magnetic interference can cause significant navigation errors.}{catD}
\dictentry{Consensus Algorithm}{A distributed algorithm that enables multiple drones in a swarm to reach agreement on a shared state, decision, or action without centralized control. Consensus algorithms ensure all agents converge on the same value through iterative message passing and local computation. They are fundamental to swarm coordination, allowing distributed decision-making even if some communication links fail.}{catD}
\dictentry{Control Link}{The radio communication channel from the pilot's transmitter to the drone's receiver, carrying control commands such as stick inputs and switch positions. A reliable control link is critical for safe flight, as loss of this link triggers failsafe procedures. Control links typically operate on 2.4 GHz or 900 MHz bands with protocols like CRSF or ELRS for low latency and high reliability.}{catD}
\dictentry{Cooperative Navigation}{A technique where multiple drones share sensor data and computational resources to navigate more effectively than any single drone could alone. By combining GPS, visual, and inertial measurements from all participating aircraft, the swarm achieves improved accuracy and robustness against individual sensor failures. This approach is especially valuable in GPS-denied environments where collaborative localization is needed.}{catD}
\dictentry{Crop Monitoring}{The aerial surveillance of agricultural fields using drone-mounted sensors to assess crop health, growth stage, and stress indicators. Drones equipped with multispectral or NDVI cameras can detect issues like nutrient deficiencies, pest infestations, and water stress before they become visible to the naked eye. This data enables precision agriculture, allowing farmers to target treatments only where needed.}{catD}
\dictentry{Crossfire}{A long-range radio control link system developed by Team Blacksheep (TBS) operating on the 868/915 MHz band. Crossfire provides extremely reliable control and telemetry at ranges exceeding 40 kilometers with very low latency. It is widely used in long-range FPV flying and professional drone operations where link reliability at great distances is essential.}{catD}
\dictentry{CRSF}{Crossfire Serial Protocol, a digital communication protocol developed by TBS for low-latency communication between the flight controller and radio transmitter. CRSF transmits channel data and telemetry bidirectionally with sub-millisecond latency, making it one of the fastest protocols available. It has become a de facto standard adopted by many RC system manufacturers beyond TBS.}{catD}
\dictentry{Cruise Speed}{The most efficient airspeed at which a drone can fly to maximize range or endurance on a given battery charge. Flying significantly above or below cruise speed increases power consumption per unit of distance traveled. For fixed-wing drones, cruise speed is typically around 60-70\% of the maximum speed, while multirotors optimize cruise speed for minimal energy per distance.}{catD}
\dictentry{Cube}{A modular flight controller system, originally developed by the Pixhawk team, featuring a cube-shaped processor module that snaps into different carrier boards. The Cube supports multiple sensor configurations and processor options, allowing easy upgrades without replacing the entire system. It is widely used in commercial and research drones for its reliability, flexibility, and extensive software support.}{catD}
\dictentry{Current Sensor}{A sensor that measures the electrical current flowing from the battery to the drone's motors and electronics in real time. This data enables the flight controller to calculate remaining battery capacity, power consumption, and estimated flight time remaining. Accurate current sensing is essential for battery management and for preventing dangerous over-discharge conditions during flight.}{catD}
\dictentry{Data Link}{A bi-directional communication channel carrying both telemetry data from the drone and commands to the drone simultaneously. Unlike a simple control link, a data link supports two-way information flow for mission updates, parameter changes, and live sensor data. Data links are critical for beyond-visual-line-of-sight operations where continuous situational awareness is required.}{catD}
\dictentry{Delivery Drone}{A drone specifically designed for the autonomous transport and delivery of packages, food, medicine, or other lightweight goods. Delivery drones incorporate GPS navigation, sense-and-avoid systems, and payload release mechanisms to complete last-mile delivery missions. They represent a rapidly growing commercial application with companies developing drone fleets for urban and rural logistics.}{catD}
\dictentry{Delivery Mechanism}{A system mounted on a delivery drone that enables controlled release of cargo at the designated location. Mechanisms range from simple servo-actuated hooks to enclosed pods with automated doors, and must ensure the payload is released safely without disturbing the drone's flight stability. The design must account for payload weight, size, and the required precision of drop-off.}{catD}
\dictentry{Demag Compensation}{A firmware feature that compensates for the demagnetization effect in brushless motors caused by back EMF during commutation. Without compensation, the motor loses torque efficiency at certain rotor positions, leading to rough performance at low throttle. Demag compensation adjusts ESC timing to counteract this effect, resulting in smoother motor operation especially during low-RPM maneuvers.}{catD}
\dictentry{Desync}{A condition where the electronic speed controller loses synchronization with the brushless motor's rotor position, causing the motor to stutter, vibrate, or stop spinning entirely. Desync typically occurs during rapid throttle changes, aggressive maneuvers, or when motor timing settings are incorrect. It can be mitigated by adjusting motor timing, reducing prop size, or using ESC firmware with better sync algorithms.}{catD}
\dictentry{Digital Twin}{A virtual replica of a physical drone system that mirrors its real-world state, behavior, and environment in a simulation or monitoring platform. Digital twins receive real-time data from sensors on the physical drone and use it to predict maintenance needs, test software updates, and simulate mission scenarios. This technology enables proactive maintenance and risk-free testing of new configurations before deployment.}{catD}
\dictentry{Direction Lock}{A firmware feature that fixes the motor rotation direction in software rather than requiring physical wire swapping. When enabled, the flight controller automatically reverses the command to any motor that is spinning in the wrong direction during initial setup. This simplifies drone assembly and maintenance by eliminating the need to resolder ESC wires when a motor is installed backwards.}{catD}
\dictentry{Directional Antenna}{An antenna that focuses its radio signal in a specific direction rather than radiating equally in all directions. Directional antennas provide greater range and signal strength in the aimed direction but require proper alignment with the target. They are used for long-range FPV video reception and control links where maximum range is more important than coverage flexibility.}{catD}
\dictentry{Disarming}{The process of disabling motor power on the drone, making it safe to handle on the ground. Disarming is typically done by a specific stick combination on the transmitter, a dedicated switch, or automatically after landing. It immediately cuts power to all motors, preventing accidental spin-up and potential injury during ground operations or maintenance.}{catD}
\dictentry{Disaster Response}{The use of drones for rapid damage assessment, search and rescue coordination, and delivery of emergency supplies following natural disasters or man-made catastrophes. Drones can survey affected areas quickly, map damage extent, locate survivors, and deliver medical supplies to inaccessible locations. Their ability to operate without ground infrastructure makes them invaluable in disaster scenarios.}{catD}
\dictentry{Distributed Control}{A swarm architecture where flight control decisions are made locally by each drone rather than by a central controller. Each agent perceives its environment, communicates with neighbors, and makes decisions based on local rules and shared information. This approach provides greater resilience since the swarm can continue operating even if individual drones or communication links fail.}{catD}
\dictentry{Diversity Receiver}{An RC receiver that uses multiple antennas, often of different types or orientations, to maintain the strongest possible control link. The receiver continuously monitors signal quality on each antenna and switches to or combines the best available signal. Diversity reception significantly reduces the chance of signal loss caused by antenna shadowing or multipath interference.}{catD}
\dictentry{DJI OSD}{The On-Screen Display system integrated into DJI's digital FPV ecosystem, providing pilots with real-time flight data overlaid on the video feed. It shows information including battery voltage, flight time, signal strength, altitude, and distance from the home point. The DJI OSD is highly integrated with DJI's ecosystem and provides seamless data visualization in their goggles and apps.}{catD}
\dictentry{Drone}{An unmanned aerial vehicle (UAV) that can be operated remotely by a pilot or fly autonomously using onboard computers and sensors. Drones range in size from tiny nano-class devices to large fixed-wing aircraft, and serve applications from recreational flying to military surveillance and commercial delivery. Modern drones integrate flight controllers, GPS, cameras, and communication systems into compact, capable platforms.}{catD}
\dictentry{Drone Registration}{The process of officially recording a drone with the national aviation authority, which assigns a unique identification number to the aircraft. Most countries require drones above a certain weight threshold to be registered, and the registration number must be displayed on the aircraft. Registration helps authorities track drone ownership, enforce safety regulations, and investigate unauthorized operations.}{catD}
\dictentry{Drone Swarm}{A group of multiple drones operating together as a coordinated entity using shared communication and control algorithms. Swarms can perform tasks collectively that would be impossible for individual drones, such as large-area search, synchronized displays, or cooperative cargo transport. Swarm technology draws from biological models like bee colonies and ant formations to achieve emergent group behavior.}{catD}
\dictentry{Dronecode Foundation}{An organization hosted by the Linux Foundation that supports and promotes open-source drone software projects including PX4, MAVLink, and QGroundControl. The foundation provides governance, funding, and infrastructure for the development of common drone standards and tools. It brings together industry partners and developers to advance the drone ecosystem through open collaboration.}{catD}
\dictentry{DroneKit}{A Python library that provides a high-level API for communicating with and controlling drones via the MAVLink protocol. DroneKit enables developers to create custom drone applications for mission planning, telemetry monitoring, and autonomous flight control without dealing with low-level protocol details. It supports both in-air and simulated drones, making it useful for testing and rapid prototyping.}{catD}
\dictentry{DShot}{A digital communication protocol between the flight controller and electronic speed controllers that provides faster, more precise throttle commands than analog protocols like PWM or Oneshot. DShot sends a fixed-length digital packet containing the throttle value and a checksum for error detection. Available in speeds from DShot150 to DShot1200, it eliminates the need for ESC calibration and provides built-in telemetry channels.}{catD}
\dictentry{DSMX}{Spektrum's proprietary radio control protocol that uses frequency-hopping spread spectrum (FHSS) for reliable, interference-resistant communication. DSMX operates on the 2.4 GHz band and provides fast frame rates with low latency for responsive drone control. It is widely used in Spektrum transmitter and receiver ecosystems with strong interference rejection in crowded RF environments.}{catD}
\dictentry{DSSS}{Direct Sequence Spread Spectrum, a modulation technique that spreads a signal over a wider bandwidth using a pseudo-random code sequence. DSSS provides strong resistance to interference and jamming because the spread signal appears as low-level noise to unauthorized receivers. It is used in some drone communication and navigation systems where signal robustness in contested environments is required.}{catD}
\dictentry{Dynamic Filter}{An adaptive notch filter in the flight controller firmware that automatically tracks and removes vibration-induced noise frequencies from sensor data. As motor RPM changes during flight, the dominant vibration frequencies shift, and dynamic filters adjust their center frequency in real time to stay aligned. This results in cleaner gyro data and smoother flight performance across the entire throttle range.}{catD}
\dictentry{EASA}{The European Union Aviation Safety Agency, responsible for establishing and enforcing aviation safety standards across EU member states. EASA sets regulations for civil drones, including operational categories, pilot licensing requirements, and aircraft certification standards. Its drone regulations implement a risk-based classification system ranging from the open category for low-risk operations to the certified category.}{catD}
\dictentry{Emergency Stop}{An immediate motor kill command that instantly cuts power to all motors, typically activated by a dedicated switch or button on the transmitter. It is used as a last resort when the drone is in a dangerous situation such as a person entanglement or imminent collision with critical infrastructure. Unlike normal disarming, emergency stop bypasses all safety checks and stops motors regardless of flight state.}{catD}
\dictentry{EmuFlight}{A fork of Betaflight firmware focused on providing exceptionally smooth and responsive flight performance through refined filtering and PID algorithms. EmuFlight incorporates innovations in gyro processing and motor output smoothing to reduce latency and improve handling feel. It is favored by freestyle pilots who prioritize smooth, locked-in flight characteristics for creative aerial maneuvers.}{catD}
\dictentry{End to End Latency}{The total time delay from when the camera captures an image to when it appears in the pilot's goggles or monitor. This includes camera processing time, video encoding, transmission delay, decoding, and display rendering. End-to-end latency directly impacts piloting responsiveness, with values under 30 milliseconds considered excellent for FPV racing and freestyle flying.}{catD}
\dictentry{Endurance}{The maximum amount of time a drone can remain in the air on a single battery charge under typical operating conditions. Endurance is determined by battery capacity, total weight, motor efficiency, and flight speed, with larger batteries providing more energy but also more weight. Endurance is a critical specification for applications like surveillance, mapping, and delivery where mission duration is paramount.}{catD}
\dictentry{Environmental Monitoring}{The use of drone-mounted sensors to track and measure environmental parameters such as air quality, water conditions, vegetation health, and wildlife activity. Drones provide a flexible, cost-effective platform for collecting environmental data at scales between ground surveys and satellite observations. This data supports conservation efforts, pollution tracking, and climate research.}{catD}
\dictentry{ESC}{An Electronic Speed Controller that regulates the rotational speed of a brushless motor by converting digital throttle commands from the flight controller into precisely timed electrical pulses. Each motor on a multirotor drone requires its own ESC to independently control thrust. Modern ESCs run specialized firmware such as BLHeli32 or AM32 that provides features like DShot support, telemetry, and active braking.}{catD}
\dictentry{ESC Calibration}{The process of calibrating electronic speed controllers to ensure they correctly interpret the full range of throttle commands from the flight controller. Calibration synchronizes the minimum and maximum pulse widths so that all motors respond uniformly to the same throttle input. Proper ESC calibration prevents inconsistent motor response, uneven thrust, and potential instability during flight.}{catD}
\dictentry{EVLOS}{Extended Visual Line of Sight, an operational mode where the drone pilot relies on trained observers positioned along the flight path to maintain visual contact. These observers relay information about the drone's position and surrounding airspace to the pilot via radio communication. EVLOS enables operations beyond a single pilot's visual range without requiring full BVLOS certification and detect-and-avoid systems.}{catD}
\dictentry{ExpressLRS}{An open-source radio control link protocol designed for extremely low latency and long-range performance. ELRS operates on both 2.4 GHz and 900 MHz bands, using LoRa modulation for its link budget and frequency hopping for interference rejection. Its open-source nature allows community-driven development, with update rates up to 1000 Hz providing near-instantaneous control response for racing and freestyle pilots.}{catD}
\dictentry{F4 Processor}{An STM32F4 series microcontroller commonly used in flight controllers, offering sufficient processing power for most flight control tasks. The F4 provides adequate computational capability for running Betaflight or INAV with standard features like OSD, logging, and multiple UART peripherals. It represents a cost-effective choice for racing drones where the latest processing power is not strictly necessary.}{catD}
\dictentry{F7 Processor}{An STM32F7 series microcontroller used in higher-end flight controllers, offering significantly more processing power than the F4. The F7 enables features like bi-directional DShot at high loop rates, advanced filtering, and more UART ports for connecting multiple peripherals simultaneously. Its additional processing headroom supports future firmware features and more complex drone configurations.}{catD}
\dictentry{FAA}{The Federal Aviation Administration, the United States government agency responsible for regulating all aspects of civil aviation, including drone operations. The FAA establishes rules for drone registration, pilot certification, airspace authorization, and operational limitations under Part 107 and other regulations. It manages the LAANC system for airspace access and enforces compliance through inspections and penalties.}{catD}
\dictentry{Failsafe}{An automatic safety action triggered by the flight controller when a critical system failure or anomaly is detected during flight. Failsafe conditions include loss of control link, low battery, compass failure, and GPS loss, each triggering an appropriate response such as return-to-home or controlled landing. Failsafe programming is essential for preventing flyaways, crashes, and unsafe conditions when the pilot cannot intervene.}{catD}
\dictentry{FHSS}{Frequency Hopping Spread Spectrum, a modulation technique where the carrier frequency rapidly switches among many discrete frequencies in a pseudorandom sequence. FHSS provides excellent resistance to interference and eavesdropping because jamming any single frequency only affects a tiny fraction of the transmitted data. It is widely used in drone radio control systems for robust communication in crowded RF environments.}{catD}
\dictentry{Fixed-Wing Drone}{A drone with rigid wings that generates lift through forward air movement, similar to a conventional airplane. Fixed-wing drones are significantly more efficient than multirotors for long-range and long-duration missions because they do not require constant motor power to stay airborne. They require a runway or launcher for takeoff and cannot hover in place, making them ideal for surveying large areas.}{catD}
\dictentry{Flight Authorization}{Official permission required to operate a drone in controlled airspace, typically issued through digital systems like the FAA's LAANC or drone zone portal. Without flight authorization, drone operations in controlled airspace near airports and heliops are illegal and dangerous. Authorization specifies the approved altitude, time window, and operational conditions for the flight.}{catD}
\dictentry{Flight Controller}{The onboard computer that processes sensor data and pilot commands to manage drone flight dynamics, stability, and autonomous behaviors. It runs firmware such as Betaflight, ArduPilot, or PX4 and interfaces with the IMU, GPS, ESCs, and radio receiver. The flight controller is the brain of the drone, executing PID control loops thousands of times per second to maintain stable flight.}{catD}
\dictentry{Flight Mode}{A specific configuration of the flight controller's control behavior that determines how the drone responds to pilot inputs and environmental conditions. Common modes include angle mode for beginners, acro mode for experienced pilots, and GPS-assisted modes like position hold and return to home. Flight modes allow the pilot to adapt the drone's behavior to different mission requirements and skill levels.}{catD}
\dictentry{Flight Termination System}{A safety system designed to end the flight in a controlled manner, typically by cutting motor power and optionally deploying a parachute. It serves as the ultimate failsafe when all other safety measures have failed or the drone has strayed into a dangerous area. Flight termination systems are often required by regulations for drones operating over people or in populated areas.}{catD}
\dictentry{Flight Time}{The maximum duration a drone can remain airborne from takeoff to landing on a single battery charge. Flight time depends on battery capacity, total aircraft weight, payload, wind conditions, and flight profile, with hover flight consuming more energy per minute than efficient forward flight. Typical consumer drones achieve 20-30 minutes of flight time, while specialized long-endurance platforms can fly for several hours.}{catD}
\dictentry{Flying Weight}{The total weight of the drone at the moment of takeoff, including the airframe, all installed equipment, battery, and any payload. Flying weight directly affects flight performance, regulatory category, battery life, and maximum payload capacity. Many aviation authorities use flying weight as the primary criterion for drone classification and applicable regulations.}{catD}
\dictentry{Foam Damping}{Foam material strategically placed between the flight controller and the drone frame to absorb high-frequency vibrations from motors and propellers. Foam damping isolates the sensitive gyroscope and accelerometer from mechanical noise that can degrade stabilization performance. The foam density and thickness are chosen to match the vibration frequency range for optimal isolation.}{catD}
\dictentry{Follow Me}{An autonomous flight mode in which the drone tracks and follows a moving subject using GPS from the pilot's phone or wearable beacon. The drone maintains a predetermined distance, altitude, and angle relative to the subject as it moves. This mode is popular for action sports filming, hiking documentation, and solo content creation where a camera operator is not available.}{catD}
\dictentry{Formation Flight}{The coordinated flight of multiple drones maintaining specific relative positions and distances throughout a mission. Formation flight requires precise synchronization of navigation, speed, and trajectory between all participating aircraft. Applications include synchronized light shows, agricultural spraying patterns, and research into aerodynamic benefits of flying in formation.}{catD}
\dictentry{FPV Camera}{A small, lightweight camera mounted on the drone that captures real-time video and transmits it wirelessly to the pilot's goggles or monitor. FPV cameras are optimized for low latency rather than image quality, with processing delays under 10 milliseconds being essential for responsive piloting. They typically use CMOS sensors with wide dynamic range to handle rapidly changing light conditions during flight.}{catD}
\dictentry{FPV Drone}{A First Person View drone piloted through a live video feed transmitted from an onboard camera to the pilot's goggles or monitor. FPV drones provide an immersive flying experience that gives the pilot a cockpit perspective, enabling precise and agile maneuvering. They are popular for racing, freestyle flying, and cinematic filmmaking, with both analog and digital video systems available.}{catD}
\dictentry{FPV Goggles}{Head-mounted displays that receive live video from the drone's FPV camera, providing the pilot with an immersive first-person flying perspective. FPV goggles vary from basic analog receivers to advanced digital systems with high resolution, wide field of view, and head tracking capabilities. They are the primary display method for FPV pilots and significantly enhance situational awareness during flight.}{catD}
\dictentry{FPV Monitor}{A ground-based screen for viewing the live FPV video feed as an alternative to goggles. Monitors are useful for spectators, trainers monitoring a student pilot, or situations where goggles are impractical. They typically include a built-in video receiver and display flight data overlaid on the video, though they provide a less immersive experience than goggles.}{catD}
\dictentry{Frame}{The structural chassis of a drone that holds all components together and provides mounting points for motors, flight controller, battery, and payload. Frame materials range from lightweight carbon fiber for racing drones to durable aluminum alloys for industrial applications. Frame geometry determines the drone's motor configuration, weight distribution, and overall flight characteristics.}{catD}
\dictentry{Frame Rate}{The number of video frames captured and transmitted per second, measured in frames per second (FPS). Standard FPV systems operate at 30 or 60 FPS, with higher frame rates providing smoother video for easier piloting. Higher frame rates require more bandwidth and can increase latency, so pilots must balance smoothness against transmission reliability.}{catD}
\dictentry{Frequency Band}{The specific radio frequency range used for a particular drone communication link, such as control, telemetry, or video. Different bands offer different tradeoffs between range, data capacity, obstacle penetration, and regulatory constraints. Common drone bands include 2.4 GHz for control, 900 MHz for long-range telemetry, and 5.8 GHz for FPV video.}{catD}
\dictentry{Gazebo}{A 3D robotics simulator developed for testing and validating drone software in a realistic physics-based virtual environment. Gazebo provides accurate simulation of sensor inputs, environmental conditions, and vehicle dynamics, enabling developers to test algorithms before deploying them on real hardware. It integrates with ROS and PX4, making it a standard tool in drone research and development.}{catD}
\dictentry{Geofence}{A virtual boundary defined in the flight controller or ground station software that restricts the drone's flight to a predetermined area. When the drone approaches or reaches the geofence boundary, it automatically takes corrective action such as stopping, hovering, or returning to the launch point. Geofences are a critical safety feature for preventing drones from entering restricted airspace or flying beyond safe operational limits.}{catD}
\dictentry{Geofence Action}{The specific response the drone executes when it reaches or violates its defined geofence boundary. Common geofence actions include hover in place, return to home, land immediately, or simply prevent further movement in the restricted direction. The appropriate action depends on the mission requirements, with return to home being the most common default for safety.}{catD}
\dictentry{Ghost}{An ultra-low latency radio control link system developed by Flywoo, designed for competitive FPV racing where minimal control delay is critical. Ghost uses LoRa modulation on the 2.4 GHz band to achieve sub-millisecond latency while maintaining excellent range and link reliability. Its lightweight receivers and compact form factor make it popular for weight-sensitive racing drone builds.}{catD}
\dictentry{Gimbal}{A motorized stabilizing mount that keeps cameras and sensors pointed in a desired direction regardless of the drone's attitude changes. Gimbals use brushless motors and inertial sensors to actively counteract the drone's pitch, roll, and yaw movements, producing smooth, stable footage. They are essential for aerial photography, inspection work, and any application requiring steady imagery from a moving platform.}{catD}
\dictentry{GLONASS}{The Global Navigation Satellite System operated by Russia, providing positioning data alongside GPS and other GNSS constellations. Many modern drone GPS modules can receive GLONASS signals simultaneously with GPS, increasing the number of visible satellites and improving positioning accuracy and reliability. Using multiple GNSS constellations provides better performance in challenging environments like urban canyons.}{catD}
\dictentry{GPS Failsafe}{An automatic action triggered by the flight controller when GPS satellite signal quality drops below acceptable levels. Common responses include switching to non-GPS flight modes, initiating a controlled landing, or alerting the pilot. GPS failsafe prevents the drone from relying on inaccurate position data that could cause erratic behavior in GPS-dependent autonomous modes.}{catD}
\dictentry{GPS Mast}{A raised mounting platform that elevates the GPS antenna above the drone's main body to improve satellite signal reception. The mast distances the GPS antenna from electromagnetic interference generated by motors, ESCs, and video transmitters. A clear sky view from the raised position also reduces signal blockage from the drone's own frame and components.}{catD}
\dictentry{GPS Module}{A receiver unit that processes signals from global navigation satellite systems to provide real-time position, velocity, and time data for drone navigation. GPS modules range from basic single-constellation receivers to advanced multi-band units supporting GPS, GLONASS, Galileo, and BeiDou simultaneously. This position data enables autonomous flight modes, geofencing, and accurate mission waypoint navigation.}{catD}
\dictentry{Gyro Calibration}{The process of calibrating the drone's gyroscope sensors to eliminate bias and ensure accurate angular velocity measurements. During calibration, the drone must remain perfectly still so the sensor can measure and record its inherent zero-rate offset. Accurate gyro calibration is essential for stable flight, as even small biases can cause the drone to drift or behave unpredictably.}{catD}
\dictentry{Gyro Sampling}{The rate at which the gyroscope sensor is read by the flight controller's processor, measured in kilohertz. Higher sampling rates provide more responsive and accurate angular velocity data for the PID control loop, with typical rates ranging from 4 kHz to 8 kHz. The sampling rate directly affects the flight controller's ability to respond to disturbances and maintain stable flight.}{catD}
\dictentry{Gyroscope}{A micro-electromechanical sensor that measures angular velocity around three orthogonal axes, detecting rotational movement of the drone. The gyroscope provides the primary data for the flight controller's stabilization algorithms, enabling it to correct for unwanted rotations and maintain the desired attitude. Combined with the accelerometer, it forms the core of the inertial measurement unit.}{catD}
\dictentry{H.264}{A widely used video codec that provides good compression efficiency for FPV video transmission. H.264 achieves a strong balance between video quality, compression ratio, and encoding speed, making it suitable for real-time video streaming. It is the most common codec in digital FPV systems and supports a wide range of resolutions and frame rates.}{catD}
\dictentry{H.265}{A more advanced video codec, also known as HEVC, that achieves approximately 50\% better compression than H.264 at equivalent visual quality. H.265 is increasingly used in digital FPV systems to deliver higher resolution video within the same bandwidth constraints. Its improved compression allows for sharper video feeds, though encoding requires more processing power and can add slight latency.}{catD}
\dictentry{H7 Processor}{An STM32H7 series microcontroller used in premium flight controllers, offering the highest processing power available for drone applications. The H7 supports features like high loop rates, advanced RPM filtering, and numerous peripheral connections for complex drone builds. Its generous processing headroom ensures the firmware can handle future feature additions without performance compromise.}{catD}
\dictentry{HALE Drone}{High Altitude Long Endurance, a class of unmanned aircraft designed to operate above 30,000 feet for periods measured in days or weeks. HALE drones use lightweight composite structures and efficient propulsion to carry sensor payloads at the edge of the stratosphere. They serve roles in surveillance, atmospheric research, and telecommunications relay, filling the gap between conventional drones and satellites.}{catD}
\dictentry{Hexacopter}{A multirotor drone with six rotors that provides greater lift capacity and motor redundancy compared to a quadcopter. If one motor fails, a hexacopter can often continue flying and land safely due to its redundant configuration. The additional motors increase weight and power consumption but provide valuable safety margins for commercial operations carrying expensive payloads.}{catD}
\dictentry{HITL}{Hardware-in-the-Loop simulation where the actual flight controller hardware runs real firmware while interfacing with a simulated drone model on a computer. HITL testing validates the interaction between real hardware and software without risking a physical aircraft. It is used for testing firmware changes, validating fail-safe behaviors, and training autopilot algorithms in realistic scenarios.}{catD}
\dictentry{Horizon Mode}{A hybrid flight mode that combines characteristics of angle mode and acro mode, providing self-leveling at moderate tilt angles while allowing flips and rolls at larger stick deflections. In horizon mode, the drone returns to level when the sticks are centered, but experienced pilots can perform acrobatic maneuvers by pushing the sticks further. It serves as a transitional mode for pilots learning to fly acro.}{catD}
\dictentry{Hover Time}{The maximum duration a drone can maintain a stationary hover at a fixed position and altitude on a single battery charge. Hover time is typically shorter than forward flight endurance because multirotors consume more energy per minute while hovering than during efficient forward flight. This specification is critical for applications like inspection and surveillance where the drone must remain in a fixed position.}{catD}
\dictentry{Hybrid Drone}{A drone that combines vertical takeoff and landing capability of a multirotor with the efficient forward flight of a fixed-wing aircraft. The drone uses rotors for vertical ascent and descent, then transitions to wing-borne flight for efficient cruise over long distances. Hybrid designs offer the operational flexibility of a multirotor with the range and endurance advantages of a fixed-wing platform.}{catD}
\dictentry{Hyperspectral Sensor}{A specialized imaging sensor that captures data across many narrow spectral bands simultaneously, far more than a standard camera or multispectral sensor. Hyperspectral imaging reveals detailed information about material composition, vegetation health, and chemical properties invisible to conventional cameras. These sensors are used in precision agriculture, mineral exploration, and environmental monitoring from drones.}{catD}
\dictentry{I2C}{Inter-Integrated Circuit, a two-wire serial communication bus used to connect low-speed peripherals to the flight controller. I2C uses a master-slave architecture with clock and data lines, allowing multiple devices to share the same bus with unique addresses. It is commonly used for connecting barometers, magnetometers, and other sensors that do not require high-speed data transfer.}{catD}
\dictentry{ICAO}{The International Civil Aviation Organization, a specialized United Nations agency that sets global standards for aviation safety, security, and environmental protection. ICAO develops framework regulations for unmanned aircraft systems that member states adapt into national legislation. Its standards ensure international consistency in drone operations, airworthiness requirements, and pilot qualifications.}{catD}
\dictentry{IMU}{An Inertial Measurement Unit that combines accelerometers and gyroscopes into a single sensor package to measure both linear acceleration and angular velocity across three axes. The IMU provides the fundamental data that the flight controller uses to determine the drone's attitude and motion in real time. Higher-end IMUs may also include magnetometers and barometers for complete inertial navigation capability.}{catD}
\dictentry{INAV}{An open-source flight controller firmware derived from Cleanflight, optimized for GPS-assisted navigation and autonomous flight capabilities. INAV focuses on features like waypoint navigation, return-to-home, position hold, and altitude hold with a user-friendly configuration interface. It is popular among long-range FPV pilots and autonomous drone builders who prioritize navigation features over racing performance.}{catD}
\dictentry{Infrastructure Inspection}{The practice of using drones to visually inspect bridges, power lines, cell towers, wind turbines, and other critical infrastructure. Drones can access hard-to-reach areas safely without scaffolding, cranes, or helicopter support, significantly reducing inspection costs and risks. High-resolution cameras and specialized sensors capture detailed imagery that engineers analyze for structural defects and maintenance needs.}{catD}
\dictentry{Inspection Drone}{A drone specifically configured or designed for close-up visual inspection of structures, equipment, and industrial facilities. Inspection drones typically feature high-resolution cameras, zoom lenses, and sometimes thermal or ultrasonic sensors for detecting subsurface defects. They reduce the need for human workers to access dangerous or elevated locations while providing thorough, repeatable documentation.}{catD}
\dictentry{IP Rating}{Ingress Protection rating, a standardized classification indicating the degree of protection a drone or component provides against dust and water intrusion. IP ratings consist of two digits, the first for dust protection and the second for water resistance, with higher numbers indicating greater protection. An IP54 rating, for example, means limited dust ingress protection and resistance to splashing water from any direction.}{catD}
\dictentry{IPX Connector}{A miniature RF connector commonly used for antenna connections on FPV video transmitters and receivers. IPX connectors, also known as U.FL or MHF connectors, provide a compact, snap-on connection suitable for the small form factors of drone electronics. They offer good RF performance up to several GHz but are delicate and intended for limited mating cycles during assembly.}{catD}
\dictentry{Kakute}{A series of flight controllers manufactured by Holybro, designed primarily for racing and freestyle FPV drones. Kakute controllers are known for their clean layout, integrated features like current sensing and OSD, and compatibility with Betaflight firmware. The series ranges from budget-friendly F4 models to high-performance F7 and H7 variants for competitive racing applications.}{catD}
\dictentry{Kiss}{A proprietary flight controller firmware and hardware ecosystem developed by Flyduino, known for its streamlined approach to drone setup and tuning. KISS simplifies the configuration process compared to open-source alternatives while providing excellent flight performance through carefully optimized defaults. The closed-source nature of KISS means it is tightly integrated with specific KISS-compatible hardware.}{catD}
\dictentry{LAANC}{Low Altitude Authorization and Notification Capability, an FAA system that provides automated, near-real-time airspace authorizations for drone flights in controlled airspace. Pilots submit flight requests through LAANC-approved apps, and the system automatically checks against airspace restrictions and issues approvals within seconds. LAANC dramatically streamlines the process of obtaining permission to fly near airports.}{catD}
\dictentry{Land Immediately}{A failsafe action where the flight controller commands the drone to land at its current location as quickly and safely as possible. This action is typically triggered by critical failures like severe battery depletion, complete loss of control link, or sensor failures that make continued flight unsafe. The drone descends vertically using available sensors to manage the landing rate and reduce impact damage.}{catD}
\dictentry{Landing Skid}{The landing gear structure on the underside of a drone that supports the aircraft when on the ground and absorbs impact during landing. Skids protect the drone's belly, camera, and payload from contact with the ground during takeoff and landing. They come in various designs from simple fixed legs to retractable systems that reduce aerodynamic drag during flight.}{catD}
\dictentry{Latency}{The delay between a pilot's control input and the drone's physical response, measured in milliseconds. Lower latency provides more immediate, responsive flight feel and is critical for FPV racing and freestyle flying where split-second reactions matter. Total system latency includes transmitter processing, radio transmission, receiver processing, flight controller computation, ESC response, and motor acceleration.}{catD}
\dictentry{LBT}{Listen Before Talk, a regulatory spectrum access protocol where a device must check if a frequency channel is clear before transmitting. LBT is required in many European and international jurisdictions for drone radio systems to prevent interference with other users. It adds a small delay before each transmission but ensures fair and coordinated use of shared radio spectrum.}{catD}
\dictentry{LED Indicator}{A light-emitting diode mounted on the drone that displays status information through color, blink pattern, or brightness changes. LED indicators show the drone's arming state, GPS lock status, battery level, and error conditions to the pilot and ground crew. They are especially useful for diagnosing issues when the drone is on the ground and for visual identification during flight.}{catD}
\dictentry{Level Calibration}{The process of calibrating the accelerometer to establish the drone's true horizontal reference level. The drone is placed on a precisely level surface during calibration so the accelerometer can record the exact sensor readings for a perfectly level attitude. Even slight errors in level calibration cause the drone to drift horizontally when in self-leveling flight modes.}{catD}
\dictentry{LHCP}{Left-Hand Circular Polarization, a type of antenna polarization where the electromagnetic wave rotates counterclockwise as it propagates. LHCP is used in FPV video systems and must match between the transmitter and receiver antenna for optimal signal strength. Mixing LHCP and RHCP antennas causes approximately 20-30 dB of signal loss, so consistent polarization throughout the video link is essential.}{catD}
\dictentry{LiDAR Payload}{A laser-based scanning system mounted on a drone for creating high-resolution 3D maps of terrain, structures, and vegetation. LiDAR sends laser pulses and measures the return time to calculate precise distances, generating point cloud data with centimeter-level accuracy. Drone-mounted LiDAR payloads are used in surveying, forestry, archaeology, and autonomous navigation applications.}{catD}
\dictentry{Linear Polarization}{A type of antenna polarization where the electromagnetic wave oscillates in a single fixed plane, either vertically or horizontally. Linear polarization is simpler to implement than circular polarization but requires careful alignment between transmitting and receiving antennas. It is commonly used in lower-cost FPV systems and telemetry links where antenna orientation remains relatively stable.}{catD}
\dictentry{LiPo Battery}{A Lithium Polymer battery, the standard rechargeable power source for modern drones due to its high energy density and ability to deliver high discharge rates. LiPo cells provide a nominal voltage of 3.7V per cell, with drone batteries typically configured in series from 2S to 12S for higher voltages. Proper LiPo care, including voltage monitoring and balanced charging, is essential for safety and longevity.}{catD}
\dictentry{Loiter}{A GPS-assisted flight mode where the drone holds its current position and altitude autonomously using satellite navigation. The pilot can make small position adjustments with the sticks while the drone returns to the loiter point when sticks are centered. This mode is useful for aerial photography, surveillance, and any situation requiring the drone to remain stationary over a specific location.}{catD}
\dictentry{Lost Link}{A condition where the control link between the pilot's transmitter and the drone's receiver is completely lost, preventing the pilot from sending commands. Lost link can result from excessive range, antenna shadowing, interference, or hardware failure. The flight controller's lost link failsafe procedure typically initiates return to home or a predetermined emergency landing to prevent a flyaway.}{catD}
\dictentry{Loudspeaker}{An audio system mounted on a drone for broadcasting public announcements, instructions, or alerts to people on the ground. Drone-mounted loudspeakers are used in emergency management, crowd control, search and rescue, and public safety operations. They allow authorities to communicate with people in inaccessible areas or hazardous environments without putting personnel at risk.}{catD}
\dictentry{Low Battery Warning}{An alert triggered when the battery voltage drops below a preset threshold, notifying the pilot that remaining flight time is limited. The warning can be delivered through the OSD, audible buzzer, transmitter vibration, or LED indicators. Early warning gives the pilot time to initiate a controlled return to home or landing before the battery reaches critically low levels.}{catD}
\dictentry{Low Pass Filter}{A signal processing filter that passes frequencies below a cutoff point and attenuates higher frequencies, commonly used to reduce gyro noise. In drone flight controllers, low pass filters remove high-frequency vibrations from motor and propeller noise before the data reaches the PID controller. Adjusting the filter cutoff frequency is a key part of tuning, balancing noise reduction against control latency.}{catD}
\dictentry{LQ}{Link Quality, a metric indicating the reliability and strength of the radio control link between the transmitter and receiver. LQ is typically expressed as a percentage, with values above 95\% considered excellent and below 80\% indicating potential problems. Unlike RSSI which measures raw signal strength, LQ accounts for packet loss and provides a more practical indication of link health.}{catD}
\dictentry{Magnetometer}{A sensor that measures the strength and direction of magnetic fields, used by the drone to determine its heading relative to magnetic north. It is essential for accurate compass heading in GPS-assisted flight modes, waypoint navigation, and return-to-home functions. Magnetometers are sensitive to electromagnetic interference from motors and power wiring, requiring careful placement and calibration.}{catD}
\dictentry{MALE Drone}{Medium Altitude Long Endurance, a class of unmanned aircraft designed to operate between 10,000 and 30,000 feet for extended periods of 24 to 48 hours. MALE drones carry surveillance, reconnaissance, and communications relay payloads for military and border security missions. They bridge the gap between smaller tactical drones and high-altitude strategic platforms, offering persistent wide-area coverage.}{catD}
\dictentry{Mapping}{The process of creating accurate orthophoto maps, 3D models, or elevation maps from overlapping images captured by a drone during systematic flight passes. Mapping requires precise GPS-tagged imagery collected with consistent overlap and ground resolution. Software processes the images using photogrammetry techniques to stitch, align, and reconstruct the mapped area into usable geospatial products.}{catD}
\dictentry{MAVLink}{Micro Air Vehicle Link, a lightweight, open-source messaging protocol designed for communication between drones and ground control stations. MAVLink uses a compact binary format with message signing for security, supporting telemetry, commands, and parameter exchanges. It is the standard protocol used by ArduPilot, PX4, and many other autopilot systems and ground stations.}{catD}
\dictentry{Mavlink 2}{The second major version of the MAVLink protocol, introducing message signing for authentication, improved extensibility, and backward compatibility with MAVLink 1. MAVLink 2 supports flexible message definitions, variable-length payloads, and a 24-bit message ID space for future expansion. It provides stronger security features essential for commercial and military drone operations.}{catD}
\dictentry{Mavlink Inspector}{A diagnostic tool that displays all MAVLink messages being received from a drone in real time, allowing developers and pilots to verify telemetry data and diagnose communication issues. It shows message types, frequencies, and content, helping identify missing, corrupted, or unexpected data streams. MAVLink inspectors are available in ground station software and standalone applications.}{catD}
\dictentry{Mavlink Protocol}{A lightweight binary communication protocol for exchanging telemetry data, commands, and status information between drones and ground stations. The protocol uses compact message frames with headers, payloads, and checksums for reliable data transmission over various communication links. It supports both one-way telemetry streaming and bidirectional command-and-control communication.}{catD}
\dictentry{MAVSDK}{A software development kit that provides a high-level API for controlling MAVLink-compatible drones using modern programming languages. MAVSDK abstracts the low-level MAVLink protocol details into convenient functions for takeoff, landing, waypoint navigation, and camera control. It supports multiple languages including C++, Python, and Swift, making drone application development more accessible.}{catD}
\dictentry{Maximum Speed}{The highest speed the drone can achieve under optimal conditions with full throttle and no payload restrictions. Maximum speed depends on motor power, propeller efficiency, aerodynamic drag, and battery voltage. This specification is important for applications like racing, pursuit operations, and time-critical missions where covering distance quickly is essential.}{catD}
\dictentry{Mesh Network}{A decentralized communication network where each drone acts as both a transmitter and relay node, creating a self-healing network topology. Mesh networks allow drones to maintain communication over large areas by relaying messages through intermediate nodes, extending range beyond point-to-point links. This architecture provides resilience against individual node failures since traffic automatically reroutes through available paths.}{catD}
\dictentry{Micro Drone}{A small drone weighing typically under 2 kilograms, designed for indoor flying, recreational use, or lightweight professional applications. Micro drones feature compact frames, smaller batteries, and less powerful motors compared to full-size drones, resulting in shorter flight times but greater portability. Many micro drones fall below regulatory weight thresholds, simplifying legal requirements for operation.}{catD}
\dictentry{Mission}{A pre-planned sequence of waypoints and actions that the drone executes autonomously without continuous pilot input. Missions define the flight path, altitude, speed, and specific actions such as taking photos or dropping payloads at designated locations. Mission-based flying is essential for surveying, mapping, and repetitive operations where consistent, systematic coverage is required.}{catD}
\dictentry{Mission Planner}{Ground station software used to plan, simulate, and upload autonomous flight missions to the drone. Mission Planner provides map-based interfaces for placing waypoints, setting altitudes, defining actions, and reviewing the planned flight path before execution. It also supports post-flight analysis by downloading and visualizing flight logs and telemetry data from completed missions.}{catD}
\dictentry{Mission Planning}{The process of designing an autonomous flight plan that specifies waypoints, altitudes, speeds, and actions for the drone to execute without pilot input. Planning includes selecting flight paths for optimal coverage, ensuring terrain clearance, and accounting for airspace restrictions and weather conditions. Well-designed missions maximize data collection efficiency while maintaining safety margins.}{catD}
\dictentry{MMCX Connector}{A push-on RF connector used for antenna connections on FPV video transmitters and receivers. MMCX connectors are slightly larger than IPX connectors and provide a more robust snap-on connection with better durability for repeated mating cycles. They offer reliable RF performance and are common on medium-sized FPV components where connector durability is important.}{catD}
\dictentry{Motor}{An electric brushless motor that converts electrical energy from the battery into rotational motion to drive the propellers and generate thrust. Drone motors are rated by size (stator diameter and height) and KV rating, which indicates RPM per volt of applied voltage. Motor selection directly impacts thrust output, efficiency, battery life, and overall flight performance.}{catD}
\dictentry{Motor Mount}{The bracket or integrated structure that attaches the motor to the drone's arm, ensuring proper alignment and secure attachment. Motor mounts must be rigid enough to transfer thrust forces to the frame without flexing, while maintaining precise motor shaft alignment with the propeller. They are typically made from aluminum or carbon fiber for optimal strength-to-weight ratio.}{catD}
\dictentry{Motor Output}{The signal sent by the flight controller to the ESC to command a specific motor speed, expressed as a percentage of maximum or a protocol-specific digital value. Motor outputs are calculated thousands of times per second by the PID controller based on sensor data and pilot commands. Precise, rapid motor output updates are essential for smooth, responsive flight performance.}{catD}
\dictentry{Motor Timing}{The timing advance of motor commutation, which determines how early the ESC fires the motor phases relative to the rotor position. Higher timing advance can increase motor speed and power output but reduces efficiency and generates more heat. Optimal motor timing varies with motor size, propeller load, and desired performance characteristics, and is adjustable in many ESC firmware.}{catD}
\dictentry{MTOM}{Maximum Takeoff Mass, the maximum weight at which a drone is certified or designed to take off, including the aircraft, fuel or battery, payload, and any external stores. MTOM is a key regulatory classification parameter that determines which operational rules and pilot certifications apply. Exceeding MTOM can compromise structural integrity, flight performance, and regulatory compliance.}{catD}
\dictentry{MTOW}{Maximum Takeoff Weight, equivalent to MTOM, specifying the maximum weight at which a drone can safely lift off. MTOW is used interchangeably with MTOM in most regulations and technical documentation. It serves as a critical threshold for drone classification, with different regulatory frameworks applying different rules based on weight categories.}{catD}
\dictentry{Multi-Agent System}{A computational framework where multiple autonomous drones or software agents coordinate their actions to achieve shared or individual objectives. Each agent operates with its own sensors, decision-making algorithms, and communication capabilities while interacting with other agents in the environment. Multi-agent systems enable complex tasks like distributed surveillance, cooperative search, and swarm manipulation.}{catD}
\dictentry{Multirotor}{A drone configuration using multiple rotors arranged symmetrically around a central body, with thrust vectoring achieved by varying individual motor speeds. Common configurations include quadcopter (four rotors), hexacopter (six), and octocopter (eight), each offering different tradeoffs in lift capacity, redundancy, and complexity. Multirotors can hover in place and maneuver in any direction, making them highly versatile.}{catD}
\dictentry{Multishot}{An older analog ESC protocol that provided faster throttle response than standard PWM by using shorter pulse widths and higher update rates. Multishot used pulse widths of 5-25 microseconds compared to PWM's 1000-2000 microseconds, resulting in quicker motor response. It has been largely superseded by digital protocols like DShot but remains compatible with many ESCs.}{catD}
\dictentry{Multispectral Sensor}{An imaging sensor that captures data in several specific spectral bands, including wavelengths beyond the visible spectrum such as near-infrared. Multispectral data reveals information about vegetation health, water quality, and soil conditions that cannot be detected with standard cameras. Drone-mounted multispectral sensors are widely used in precision agriculture and environmental monitoring.}{catD}
\dictentry{Nano Drone}{A very small drone weighing typically under 250 grams, designed for indoor flying, educational use, or recreational enjoyment. Nano drones are lightweight enough to avoid many regulatory requirements in some countries and are safe to fly in confined spaces. Despite their small size, modern nano drones can incorporate features like altitude hold, headless mode, and beginner-friendly stabilization.}{catD}
\dictentry{No-Fly Zone}{A restricted geographic area where drone flight is prohibited by law or regulation, typically around airports, military installations, government buildings, and critical infrastructure. No-fly zones are enforced through geofencing in consumer drone firmware and through legal penalties for violations. Some no-fly zones are temporary, such as during major events or emergency operations.}{catD}
\dictentry{NOTAM}{Notice to Air Missions, an official advisory issued to alert pilots about important information regarding airspace, facilities, or procedures that may affect a flight. NOTAMs for drones can indicate temporary flight restrictions, hazard areas, or changes to airspace availability. Drone pilots must check for relevant NOTAMs before flying, especially near airports or during large public events.}{catD}
\dictentry{Notch Filter}{A signal processing filter that removes a very narrow band of frequencies while passing all others, used to target specific vibration frequencies. In drone flight controllers, notch filters are tuned to eliminate the dominant vibration frequency caused by motors and propellers. Unlike broader low-pass filters, notch filters reduce noise with minimal impact on control response and latency.}{catD}
\dictentry{Obstacle Detection}{The capability of a drone to sense obstacles in its flight path using onboard sensors such as cameras, LiDAR, ultrasonic rangefinders, or time-of-flight sensors. Detected obstacles trigger automatic avoidance responses like braking, hovering, or rerouting around the obstruction. Obstacle detection is critical for autonomous operations, indoor flying, and flights near structures or people.}{catD}
\dictentry{Octocopter}{A multirotor drone with eight rotors that provides the maximum lift capacity and motor redundancy of any common multirotor configuration. An octocopter can sustain the loss of up to two motors and still maintain controlled flight, making it the safest choice for carrying expensive cinema cameras or operating over people. The eight motors distribute the load, reducing stress on individual components.}{catD}
\dictentry{Omnidirectional Antenna}{An antenna that radiates and receives radio signals equally in all directions around its axis. Omnidirectional antennas provide consistent coverage regardless of the drone's orientation relative to the ground station, making them ideal for control links and telemetry. They offer less range than directional antennas in any single direction but provide more reliable connectivity during dynamic flight maneuvers.}{catD}
\dictentry{Oneshot}{An early high-frequency analog ESC protocol that improved throttle response over standard PWM by using shorter pulse widths of 125-250 microseconds. Oneshot allowed faster motor updates, resulting in more responsive flight characteristics compared to PWM. It was a significant improvement for racing drones but has since been largely replaced by digital protocols like DShot.}{catD}
\dictentry{Operator License}{A certification required for individuals who operate drones commercially, demonstrating knowledge of aviation regulations, safety procedures, and operational competency. Operator licenses vary by country, with the FAA's Part 107 certificate in the US and the A2 Certificate of Competency in the EU being common examples. Licensed operators can perform commercial work such as aerial surveying, photography, and inspection.}{catD}
\dictentry{Orbit Mode}{An autonomous flight mode where the drone circles around a designated point of interest at a specified radius, altitude, and speed. The drone continuously adjusts its heading to face the center point while maintaining a consistent orbit, ideal for surveillance, filming, or inspection of a structure. Orbit parameters such as radius, altitude, and direction can be adjusted during the orbit.}{catD}
\dictentry{OSD}{On-Screen Display, a system that overlays flight data and telemetry information directly onto the FPV video feed in the pilot's goggles or monitor. The OSD displays critical information including battery voltage, current draw, flight time, GPS coordinates, altitude, distance, and signal strength. OSD configuration is customizable through firmware settings, allowing pilots to prioritize the information most relevant to their flying style.}{catD}
\dictentry{Parachute System}{A recovery device that deploys a parachute to slow the drone's descent in the event of a critical system failure. Parachute systems are designed to open rapidly and reduce the drone's terminal velocity to minimize ground impact energy. They are often required by regulations for flights over populated areas and represent a final safety measure to protect people and property on the ground.}{catD}
\dictentry{Parameter}{A configurable setting stored in the flight controller's memory that controls a specific aspect of drone behavior, such as PID gains, flight mode options, or failsafe actions. Parameters can be modified through ground station software to customize the drone's performance without changing the firmware code. Most flight controllers store hundreds of parameters, each with a default value that can be adjusted.}{catD}
\dictentry{Parcel Delivery}{The autonomous transport and drop-off of packages by drones, representing a rapidly developing commercial application of unmanned aerial technology. Parcel delivery drones incorporate GPS navigation, sense-and-avoid systems, and automated payload release mechanisms for hands-free operation. Major logistics companies are developing drone delivery networks to provide faster, more efficient last-mile delivery services.}{catD}
\dictentry{Part 107}{The FAA regulations governing small unmanned aircraft systems (under 55 pounds) used for commercial operations in the United States. Part 107 establishes rules for pilot certification, airspace restrictions, altitude limits, visual line-of-sight requirements, and operational limitations. It provides the regulatory framework that enables commercial drone applications like photography, surveying, and inspection.}{catD}
\dictentry{Patch Antenna}{A flat, directional antenna that focuses radio energy in a specific forward-facing direction, providing higher gain than omnidirectional antennas. Patch antennas are commonly used for FPV video reception where the pilot knows the general direction of the drone and wants maximum signal strength. They offer a good compromise between directionality and size, fitting easily on goggles or ground stations.}{catD}
\dictentry{Payload Capacity}{The maximum additional weight a drone can carry beyond its own airframe weight while maintaining safe flight characteristics. Payload capacity is determined by the drone's motor thrust, structural strength, battery power, and regulatory weight limits. Exceeding payload capacity reduces flight time, agility, and safety margins, and may violate regulatory weight thresholds.}{catD}
\dictentry{PDB}{Power Distribution Board, a circuit board that receives power from the battery and distributes it to multiple ESCs and other high-current components. The PDB often includes voltage regulators for powering the flight controller and accessories, as well as current and voltage sensing for battery monitoring. Some modern PDBs integrate with ESCs into all-in-one boards to save weight and simplify wiring.}{catD}
\dictentry{Photogrammetry}{The science of creating accurate 3D models and measurements from overlapping photographs taken from multiple angles. Drone photogrammetry uses GPS-tagged aerial images with sufficient overlap to reconstruct terrain, buildings, and objects as detailed 3D point clouds or textured meshes. Software like Pix4D and OpenDroneMap processes the images using structure-from-motion algorithms to generate survey-grade outputs.}{catD}
\dictentry{PID Loop Frequency}{The rate at which the flight controller executes its PID control loop, measured in kilohertz. Higher loop frequencies enable more responsive corrections and smoother flight, with typical values ranging from 4 kHz to 8 kHz in modern firmware. The loop frequency must be supported by adequate sensor sampling rates and processor capability to avoid timing overruns.}{catD}
\dictentry{PID Tuning}{The process of adjusting the proportional, integral, and derivative gains of the flight controller's PID algorithm to optimize flight performance. Proper PID tuning ensures the drone responds accurately to pilot commands without oscillating, overshooting, or feeling sluggish. Tuning can be done manually through flight testing or automatically using autotune features, with optimal values varying for each drone configuration.}{catD}
\dictentry{Pitch Propeller}{A propeller specifically designed for clockwise rotation, identified by blade markings or numbering. In a multirotor setup, alternating motors use pitch (clockwise) and roll (counter-clockwise) propellers to cancel rotational torque. Using the wrong propeller direction on a motor causes the drone to spin uncontrollably and flip on takeoff.}{catD}
\dictentry{Pixhawk}{A widely used open-source flight controller hardware platform that supports ArduPilot and PX4 firmware. Pixhawk boards feature comprehensive sensor suites, multiple UART ports, and robust connectors for professional and research drone applications. The Pixhawk ecosystem includes various form factors and versions, from the original full-size board to mini and cube variants for different drone sizes.}{catD}
\dictentry{Position Hold}{A GPS-assisted flight mode where the drone maintains its current horizontal position and altitude using satellite navigation data. The pilot can reposition the drone with stick inputs, and it will hold the new position when the sticks are released. Position hold is essential for aerial photography, surveillance, and inspections where the drone must remain steady over a specific location.}{catD}
\dictentry{Power Distribution Board}{A board that receives high-current power from the battery and distributes it to each ESC and high-power component on the drone. It often includes integrated voltage regulators to provide stable 5V or 12V supplies for the flight controller and accessories. Modern PDBs may incorporate current sensing and OSD functionality for comprehensive battery monitoring.}{catD}
\dictentry{Power Module}{A module placed between the battery and the PDB that measures the battery's voltage and current draw in real time. The data is sent to the flight controller for OSD display, battery failsafe calculations, and remaining capacity estimation. Power modules are essential for accurate battery management and preventing dangerous over-discharge during flight.}{catD}
\dictentry{Pre-Arm Check}{A series of automated safety verifications that the flight controller performs before allowing the motors to be armed. Pre-arm checks include verifying GPS fix quality, compass calibration, accelerometer status, receiver signal, and battery voltage. If any check fails, the flight controller prevents arming and displays an error message, preventing potentially dangerous takeoffs.}{catD}
\dictentry{Precision Agriculture}{The use of drone technology for data-driven crop management, including targeted application of water, fertilizers, and pesticides based on aerial sensor data. Drones equipped with multispectral sensors identify areas of crop stress, disease, or nutrient deficiency that are invisible to the naked eye. This approach reduces chemical usage, lowers costs, and increases crop yields compared to blanket field treatment.}{catD}
\dictentry{Propeller}{A rotating airfoil blade that generates thrust by accelerating air downward, propelling the drone forward or providing lift. Propellers are characterized by their diameter, pitch, and blade count, with different combinations optimized for speed, efficiency, or quiet operation. Correct propeller selection is critical for flight performance, as mismatched propellers cause inefficient thrust, vibrations, or motor overload.}{catD}
\dictentry{Propeller Guard}{A protective ring or frame around the propeller that prevents the spinning blades from contacting people, objects, or the ground during minor collisions. Propeller guards reduce the risk of injury and property damage, making them essential for indoor flying and operations near people. They add weight and slightly reduce aerodynamic efficiency but provide valuable safety protection.}{catD}
\dictentry{PWM}{Pulse Width Modulation, an analog signaling method where information is encoded in the width of electrical pulses. PWM is the traditional protocol for controlling servos and ESCs, with pulse widths typically ranging from 1000 to 2000 microseconds. While simple and widely compatible, PWM has been largely replaced by digital protocols like DShot that offer faster response and error checking.}{catD}
\dictentry{PX4}{An open-source flight control software for drones that provides advanced autopilot capabilities including autonomous navigation, mission planning, and multi-vehicle coordination. PX4 runs on various flight controller hardware and is backed by the Dronecode Foundation with contributions from academic and industry partners. It is widely used in research, commercial drone development, and the academic community.}{catD}
\dictentry{PX4 Community}{The global network of developers, researchers, and companies that contribute to the development and improvement of the PX4 open-source autopilot software. The community operates through GitHub, forums, and regular developer meetings, coordinating efforts to advance drone autonomy and capabilities. It includes major drone manufacturers, universities, and individual contributors from around the world.}{catD}
\dictentry{QGroundControl}{An open-source ground control station application used for configuring, planning, and monitoring drone missions. QGroundControl supports both ArduPilot and PX4 firmware, providing a unified interface for parameter configuration, mission planning, and real-time telemetry display. It runs on Windows, macOS, Linux, Android, and iOS, making it one of the most cross-platform ground station options available.}{catD}
\dictentry{Quadcopter}{A multirotor drone with four rotors arranged in a cross or X configuration, the most common drone design for both consumer and commercial applications. The four motors provide sufficient lift for cameras and payloads while keeping the design simple, lightweight, and easy to maintain. Quadcopters achieve yaw control by spinning diagonal motor pairs in opposite directions to cancel torque.}{catD}
\dictentry{Raceflight}{A specialized flight controller firmware optimized for competitive FPV drone racing, prioritizing minimal latency and maximum responsiveness. Raceflight was designed specifically for the racing community with features like extremely fast PID loop rates and simplified configuration. It has since been largely absorbed into the broader FPV firmware ecosystem but contributed innovations in low-latency control.}{catD}
\dictentry{Racing Drone}{A high-speed drone specifically built for competitive FPV racing, featuring lightweight frames, powerful motors, and minimal aerodynamic drag. Racing drones can exceed 100 mph and are designed for rapid acceleration, sharp turns, and precise handling through race courses. They sacrifice camera quality and flight time for maximum speed, agility, and crash durability.}{catD}
\dictentry{Radio Calibration}{The process of configuring the flight controller to correctly interpret the full range of signals from the pilot's radio transmitter. Calibration sets the channel endpoints, center points, and direction to ensure that stick movements translate accurately to drone commands. Incorrect radio calibration causes asymmetric control response, reversed channels, or incomplete travel range.}{catD}
\dictentry{Radio Link}{The communication channel between the pilot's radio transmitter and the drone's onboard receiver, carrying control commands, telemetry data, and sometimes audio. The radio link's reliability and range depend on transmitter power, antenna quality, frequency band, and environmental conditions. Maintaining a strong, low-latency radio link is essential for safe and responsive drone operation.}{catD}
\dictentry{Range}{The maximum distance the drone can fly from its controller while maintaining a reliable control and video link. Range is limited by transmitter power, antenna gain, receiver sensitivity, and environmental obstacles, with different protocols offering vastly different maximum ranges. Long-range systems like Crossfire or ELRS on 900 MHz can achieve ranges exceeding 40 kilometers under ideal conditions.}{catD}
\dictentry{Rate Mode}{A flight mode, also known as acro or manual mode, where the drone's angular rotation rate is directly proportional to stick deflection. The drone does not self-level, requiring continuous pilot input to maintain a desired attitude. Rate mode provides maximum control authority and is preferred by experienced pilots for aerobatics, freestyle flying, and racing.}{catD}
\dictentry{RC Failsafe}{An automatic action triggered by the flight controller when the radio control link between the transmitter and receiver is lost. Common RC failsafe actions include return to home, hover in place, land immediately, or continue on the last known course. The appropriate fail-safe action should be configured before every flight to ensure the drone behaves safely when the pilot loses control.}{catD}
\dictentry{Receiver}{The radio receiver mounted on the drone that captures control signals from the pilot's transmitter and forwards them to the flight controller. Receivers come in various protocols including CRSF, SBUS, and DSMX, each with different latency, range, and channel count characteristics. The receiver's antenna placement and orientation significantly affect signal quality and control link reliability.}{catD}
\dictentry{Reconnaissance Drone}{A drone designed for intelligence gathering, surveillance, and reconnaissance missions, typically equipped with high-resolution cameras and sensor payloads. Reconnaissance drones provide real-time situational awareness to military commanders without risking human pilots in hostile territory. They range from small hand-launched tactical drones to large medium-altitude platforms with multi-hour endurance.}{catD}
\dictentry{Remote ID}{A system that broadcasts the drone's identity, location, altitude, and operator information over radio or internet to enable identification and tracking. Remote ID is required by many aviation authorities as a safety measure to allow authorities and the public to identify drones in their airspace. It functions as a digital license plate for drones, supporting accountability and security.}{catD}
\dictentry{Remote Pilot Certificate}{A certification demonstrating that a drone pilot has passed the required knowledge examination and meets the competency standards for commercial drone operations. The certificate is issued by the national aviation authority, such as the FAA's Remote Pilot Certificate under Part 107 in the United States. It authorizes the holder to conduct commercial drone operations within the specified regulatory framework.}{catD}
\dictentry{Resolution}{The pixel dimensions of a video feed, commonly expressed as 720p, 1080p, or 4K, determining the level of detail visible in the FPV image. Higher resolutions provide sharper video for better situational awareness and recording quality but require more bandwidth and can increase transmission latency. FPV pilots typically choose resolution based on their priorities, with racing pilots favoring lower latency over image quality.}{catD}
\dictentry{Resource Remapping}{The process of reassigning the flight controller's internal pin functions to different physical output pads, allowing flexible hardware configurations. Resource remapping enables builders to use any motor output pad for any motor, and assign UART ports to different peripherals regardless of the board's default layout. This feature is essential for troubleshooting hardware issues and custom drone builds.}{catD}
\dictentry{Return to Home}{An autonomous flight mode that navigates the drone from its current position back to the GPS-recorded takeoff point and initiates a landing. RTH can be activated manually by the pilot or automatically as a failsafe response to low battery, lost link, or geofence violation. The drone typically climbs to a preset safe altitude before navigating home to avoid obstacles along the return path.}{catD}
\dictentry{Return to Launch}{A failsafe action that commands the drone to autonomously navigate back to its original takeoff point and land. Similar to return to home, return to launch is specifically triggered as an automatic safety response rather than a manually selected flight mode. It serves as the primary failsafe mechanism for preventing flyaways when the pilot loses control or awareness of the drone.}{catD}
\dictentry{RHCP}{Right-Hand Circular Polarization, a type of antenna polarization where the electromagnetic wave rotates clockwise as it propagates. RHCP is widely used in FPV video systems and must match between the transmitter and receiver antennas for optimal signal reception. Most FPV systems use either RHCP or LHCP consistently, and mixing polarization types causes severe signal degradation.}{catD}
\dictentry{ROS}{Robot Operating System, a flexible framework of libraries and tools for building robot software, widely used in drone development and research. ROS provides hardware abstraction, device drivers, inter-process communication, and package management for developing complex autonomous behaviors. It is not an operating system itself but a middleware layer that simplifies the development of sophisticated drone applications.}{catD}
\dictentry{RP-SMA}{Reverse Polarity SubMiniature version A, a threaded RF connector variant where the center pin and socket are swapped compared to standard SMA. RP-SMA connectors are commonly used for antenna connections on FPV video equipment and wireless devices. Despite the reverse polarity designation, RP-SMA connectors are physically compatible with standard SMA threads but require matching gender.}{catD}
\dictentry{RPAS}{Remotely Piloted Aircraft System, the formal ICAO term encompassing the entire system of an unmanned aircraft, its control station, data links, and support equipment. RPAS emphasizes that the pilot operates the aircraft remotely rather than being physically on board. This terminology is used in international regulations and distinguishes professional drone operations from toy-grade remote-controlled aircraft.}{catD}
\dictentry{RPM Filtering}{An advanced filtering technique that uses real-time motor RPM data to calculate and track vibration frequencies for dynamic notch filtering. As motor speed changes during flight, RPM filtering adjusts the notch filter positions in real time to always target the actual vibration frequencies. This provides significantly better vibration rejection than static filters across the entire throttle range.}{catD}
\dictentry{RSSI}{Received Signal Strength Indicator, a measurement of the radio signal strength at the receiver expressed in decibels or as a percentage. RSSI provides a general indication of link quality but does not account for packet loss, making it less reliable than Link Quality for assessing actual communication health. Declining RSSI warns of potential signal loss as the drone flies further from the transmitter.}{catD}
\dictentry{RTCA}{A private, not-for-profit organization that develops consensus-based standards for the aviation industry, including standards for unmanned aircraft systems. RTCA Special Committees produce DO-series documents that define minimum operational performance standards for drone equipment and systems. Its standards are referenced by aviation authorities worldwide for drone airworthiness and operational certification.}{catD}
\dictentry{RTH}{Return to Home, an abbreviated term for the autonomous flight mode that brings the drone back to its recorded takeoff point. RTH is one of the most important safety features in modern drones, serving as the primary failsafe response for many emergency conditions. The abbreviation is commonly used in flight controller firmware, OSD displays, and pilot communications.}{catD}
\dictentry{S Rating}{The number of cells connected in series in a LiPo battery, which determines the battery's nominal and maximum voltage. A 4S battery has four cells in series for a nominal voltage of 14.8V, while a 6S battery provides 22.2V. Higher S ratings deliver more voltage and power to the motors, enabling faster flight speeds and greater thrust, but require appropriately rated ESCs and motors.}{catD}
\dictentry{Safety Switch}{A physical switch that must be activated before the drone can be armed, providing an additional layer of protection against accidental motor start. The safety switch is typically mounted on the drone's body and requires a deliberate press before arming is permitted. It prevents unintended motor spin-up during transport, handling, or maintenance when the transmitter might accidentally command arming.}{catD}
\dictentry{SBUS}{A serial bus protocol developed by Futaba for RC receivers that transmits up to 18 channels of digital data over a single wire. SBUS uses inverted serial communication at 100,000 baud with 16-bit resolution per channel for precise control signal transmission. It has been widely adopted across the drone industry and is supported by most modern flight controllers.}{catD}
\dictentry{Scientific Instrument}{A specialized sensor payload mounted on a drone for conducting scientific measurements and research, such as atmospheric sampling, radiation detection, or environmental analysis. These instruments are calibrated for precision and accuracy, enabling researchers to collect data at scales between ground-based measurements and satellite observations. Drone-mounted instruments provide flexibility for targeted, repeatable measurements in remote or hazardous locations.}{catD}
\dictentry{SD Card Logging}{The process of recording flight data, telemetry, and sensor readings to an onboard SD card for post-flight analysis and debugging. SD card logs contain detailed information about motor outputs, pilot inputs, sensor data, and GPS tracks that can be analyzed with specialized tools. Comprehensive logging is essential for tuning PID controllers, diagnosing issues, and documenting flight performance.}{catD}
\dictentry{Search and Rescue}{The use of drones to locate missing persons, assess disaster damage, and coordinate rescue operations in challenging terrain or conditions. Drones equipped with thermal cameras, powerful lights, and loudspeakers can search large areas far more quickly than ground teams. They provide real-time aerial views to rescue coordinators and can deliver first aid supplies to inaccessible locations before ground teams arrive.}{catD}
\dictentry{Search and Rescue Drone}{A drone specifically equipped with thermal cameras, searchlights, and communication systems for locating missing persons and supporting rescue operations. SAR drones can detect heat signatures from humans even in darkness or dense vegetation, dramatically reducing search time compared to ground-based methods. They often carry loudspeakers for two-way communication with found individuals while rescue teams are en route.}{catD}
\dictentry{Serial Port}{A communication interface on the flight controller used to connect peripherals like GPS modules, receivers, telemetry radios, and other sensors. Serial ports use UART communication at configurable baud rates and are assigned specific functions through the firmware configuration. Modern flight controllers include multiple serial ports to support the various peripherals required by complex drone builds.}{catD}
\dictentry{Service Ceiling}{The maximum altitude at which a drone can maintain steady, controlled flight under standard atmospheric conditions. Service ceiling is limited by decreasing air density, which reduces both aerodynamic lift and motor cooling efficiency. Beyond the service ceiling, the drone may experience reduced control authority, loss of lift, and increased power consumption that compromise safe flight.}{catD}
\dictentry{SITL}{Software-in-the-Loop simulation where the entire autopilot software stack runs on a computer without any physical hardware, simulating all sensor inputs and vehicle dynamics in software. SITL allows developers to test and debug flight controller code rapidly without risking a real aircraft. It is invaluable for developing new features, testing failsafe behaviors, and validating mission plans before field testing.}{catD}
\dictentry{SMA Connector}{A SubMiniature version A threaded RF connector commonly used for antenna connections on FPV video transmitters, receivers, and radio equipment. SMA connectors provide a secure, vibration-resistant connection with reliable RF performance across a wide frequency range. They come in standard and reverse polarity variants, and proper mating is essential to avoid damage to the connector center pin.}{catD}
\dictentry{SNR}{Signal-to-Noise Ratio, a measurement comparing the strength of the desired signal to the background noise level, expressed in decibels. Higher SNR values indicate cleaner, more reliable communication with less data corruption from noise. In drone applications, SNR is important for assessing video link quality, radio control reliability, and sensor measurement accuracy.}{catD}
\dictentry{Solar Drone}{A drone powered by integrated solar panels that convert sunlight into electrical energy to supplement or replace battery power during flight. Solar drones can achieve dramatically extended flight endurance by harvesting energy during daylight hours, with some designs capable of multi-day flights. They are used for high-altitude reconnaissance, atmospheric monitoring, and telecommunications relay.}{catD}
\dictentry{SPI}{Serial Peripheral Interface, a high-speed synchronous communication bus used to connect sensors, flash memory, and other peripherals to the flight controller. SPI provides faster data transfer rates than I2C, making it suitable for connecting gyros and accelerometers that require high-frequency data sampling. The SPI bus uses separate chip-select lines to communicate with multiple devices independently.}{catD}
\dictentry{Spotlight}{A high-intensity light mounted on a drone for illuminating areas during nighttime operations, search and rescue, or surveillance missions. Drone-mounted spotlights can be controlled remotely to aim at specific targets while the drone hovers or flies. They extend operational capability into darkness and are essential for nighttime inspection, security, and emergency response operations.}{catD}
\dictentry{Sprayer}{A liquid dispensing system mounted on agricultural drones for applying pesticides, herbicides, fertilizers, or other liquid treatments to crops. Sprayer systems include tanks, pumps, nozzles, and flow control electronics that ensure precise, even application rates during automated flight passes. Modern sprayers use GPS-guided flight paths and variable rate technology to optimize chemical usage.}{catD}
\dictentry{Submersible Drone}{An unmanned underwater vehicle designed for marine inspection, research, and exploration beneath the water's surface. Submersible drones use thrusters for propulsion and depth-rated cameras and sensors for underwater imaging and data collection. They are used for hull inspection, aquaculture monitoring, underwater archaeology, and search operations in environments too dangerous for human divers.}{catD}
\dictentry{Surveying}{The practice of measuring and mapping land features, elevation, and boundaries using drone-mounted cameras and sensors. Drone surveying combines GPS-tagged aerial photography with photogrammetry software to create accurate maps, 3D models, and elevation profiles of the terrain. It is significantly faster and more cost-effective than traditional ground surveying methods for large areas.}{catD}
\dictentry{Surveying Drone}{A drone equipped with high-resolution cameras, RTK GPS, and photogrammetry software for accurate mapping and land surveying applications. Surveying drones capture systematic aerial imagery with precise positioning data to generate survey-grade maps and 3D models. They are widely used in construction, mining, agriculture, and land management for measuring volumes, tracking progress, and documenting site conditions.}{catD}
\dictentry{Swarm Coordination}{The management and synchronization of multiple drones' actions to achieve a collective objective through shared communication and control algorithms. Coordination algorithms distribute tasks, resolve conflicts, and maintain formation integrity while allowing individual drones to respond to local conditions. Effective swarm coordination enables complex group behaviors like synchronized displays, cooperative search, and distributed sensing.}{catD}
\dictentry{Swarm Drone}{A drone specifically designed or configured to operate as part of a coordinated group, featuring compact size, reliable communication, and cooperative navigation capabilities. Swarm drones are typically lightweight and inexpensive enough to deploy in large numbers, with each unit contributing to the collective mission. They incorporate inter-drone communication and distributed decision-making algorithms.}{catD}
\dictentry{Swarm Intelligence}{Collective behavior and problem-solving capability that emerges from the interactions of many simple agents following local rules, inspired by biological systems like ant colonies and bird flocks. Swarm intelligence enables drone groups to achieve complex objectives through decentralized coordination without a central controller. This approach provides scalability, resilience, and adaptability that centralized systems cannot match.}{catD}
\dictentry{Tactical Drone}{A mid-range military drone designed for battlefield surveillance, target acquisition, and operational support at the brigade or battalion level. Tactical drones bridge the gap between small hand-launched systems and large strategic platforms, providing persistent surveillance and real-time intelligence. They are designed for rapid deployment, harsh environments, and integration with military command and control systems.}{catD}
\dictentry{Telemetry}{The real-time transmission of flight data from the drone to the ground station, providing the pilot and operators with information about the aircraft's status, position, and health. Telemetry data includes battery voltage, GPS coordinates, altitude, speed, temperature, and sensor readings. It is transmitted via the radio control link or a separate telemetry radio for monitoring and logging during flight.}{catD}
\dictentry{Tethered Drone}{A drone connected to a ground-based power supply and communication link via a physical cable, enabling unlimited flight endurance without battery constraints. Tethered drones can remain airborne for days or weeks, providing persistent surveillance, communications relay, or lighting over a fixed area. The tether limits mobility and range but eliminates the primary limitation of battery-powered flight.}{catD}
\dictentry{Thermal Camera}{A camera that detects infrared radiation emitted by objects, producing images based on temperature differences rather than visible light. Thermal cameras on drones can detect heat signatures from people, animals, electrical faults, and insulation gaps even in complete darkness. They are essential tools for search and rescue, building inspection, wildlife monitoring, and firefighting operations.}{catD}
\dictentry{Throttle Calibration}{The process of calibrating the flight controller's throttle channel to ensure the correct range of throttle input is recognized. Calibration sets the minimum and maximum pulse widths that correspond to zero and full throttle on the transmitter. Proper throttle calibration ensures consistent motor response and prevents issues like inability to disarm or unexpected full-throttle commands.}{catD}
\dictentry{Throttle Limit}{A firmware setting that caps the maximum throttle output to a percentage below 100\%, reducing the drone's maximum thrust and speed. Throttle limits are useful for training beginners, flying in confined spaces, or extending battery life by preventing aggressive power draws. The limit can often be assigned to a transmitter switch for on-the-fly adjustment.}{catD}
\dictentry{Throttle Scale}{A firmware parameter that scales the throttle input curve to provide a more linear relationship between stick position and motor output. Without scaling, the throttle response may feel non-linear due to the relationship between voltage, RPM, and thrust. Throttle scaling makes the drone feel more predictable and easier to control, especially at lower throttle settings.}{catD}
\dictentry{Tracer}{A low-latency radio control link system from Team Blacksheep (TBS) designed for racing and freestyle FPV flying where minimal delay is critical. Tracer operates on the 2.4 GHz band and provides sub-millisecond latency with reliable link quality for competitive flying. It serves as a faster alternative to Crossfire for pilots who prioritize control responsiveness over maximum range.}{catD}
\dictentry{Transmitter}{The handheld radio controller used by the pilot to send flight commands to the drone, featuring control sticks, switches, and often a built-in display. Modern transmitters support multiple protocols and can bind with receivers from various manufacturers through firmware flexibility. Transmitter quality directly affects the pilot's control precision and the reliability of the command link to the drone.}{catD}
\dictentry{UART}{Universal Asynchronous Receiver-Transmitter, a serial communication interface used on flight controllers to connect peripherals like GPS modules, telemetry radios, and receivers. UART ports communicate at configurable baud rates and support bidirectional data exchange with connected devices. Most flight controllers include multiple UARTs to accommodate the various peripherals required by complex drone builds.}{catD}
\dictentry{UAS}{Unmanned Aircraft System, the complete system comprising the unmanned aircraft, ground control station, data links, and all supporting equipment needed for operation. UAS is the formal term preferred by aviation authorities because it encompasses the entire operational system, not just the aircraft itself. This terminology emphasizes that drone operations require integrated systems rather than just a flying vehicle.}{catD}
\dictentry{UAV}{Unmanned Aerial Vehicle, the aircraft component of an unmanned aircraft system, referring specifically to the flying vehicle without its ground control and support equipment. UAV is the most widely recognized term for drones in both technical and general contexts. While sometimes used interchangeably with drone, UAV is the more formal designation used in regulatory and technical documentation.}{catD}
\dictentry{UFL Connector}{A tiny surface-mount RF connector used for antenna connections on compact FPV video transmitters and receivers. UFL connectors provide the smallest form factor of any common RF connector, making them suitable for weight-critical applications where every gram matters. They have limited durability for repeated mating cycles and are best suited for permanent installations during assembly.}{catD}
\dictentry{UTM}{Unmanned Aircraft System Traffic Management, a system for managing drone traffic in low-altitude airspace to prevent collisions and ensure safe coexistence with manned aviation. UTM provides services including airspace authorization, flight planning, real-time tracking, and separation assurance for drone operations. It is the drone equivalent of traditional air traffic control, designed for the high density and autonomous nature of unmanned operations.}{catD}
\dictentry{Vibration Damping}{The practice of isolating the flight controller and sensors from mechanical vibrations produced by the drone's motors and propellers. Damping uses rubber grommets, foam, silicone, or specialized mounting systems to absorb high-frequency oscillations before they reach the IMU. Effective vibration damping is critical for stable flight, as excessive vibration causes sensor noise that degrades stabilization performance.}{catD}
\dictentry{Video Link}{The radio communication channel that transmits live video from the drone's onboard camera to the pilot's ground-based display, typically goggles or a monitor. Video links can be analog for low latency or digital for higher image quality, with different systems offering varying ranges, resolutions, and transmission speeds. The video link is the pilot's eyes in the air and is essential for FPV operations.}{catD}
\dictentry{Video Receiver}{A ground-based unit that receives the wireless video signal transmitted from the drone's video transmitter and outputs it to goggles or a monitor. Video receivers include antenna inputs, demodulation circuitry, and often diversity switching between multiple antennas for optimal signal quality. The receiver's sensitivity and noise figure directly affect the range and clarity of the FPV video feed.}{catD}
\dictentry{VLOS}{Visual Line of Sight, an operational requirement where the drone pilot maintains direct, unaided visual contact with the drone at all times during flight. VLOS operations ensure the pilot can see the drone's attitude, detect obstacles, and avoid other aircraft without relying on cameras or instruments. Most basic drone regulations require VLOS unless specific waivers or certifications for BVLOS operations are obtained.}{catD}
\dictentry{Voltage Sensor}{A sensor that measures the drone battery's real-time voltage level and reports it to the flight controller for display on the OSD and failsafe calculations. Voltage sensing enables low-battery warnings and automatic return-to-home triggers to prevent deep discharge that could damage the battery or cause a crash. Accurate voltage sensing is essential for battery health monitoring and flight safety.}{catD}
\dictentry{VTOL Drone}{A drone capable of Vertical Takeoff and Landing like a multirotor, combined with the efficient forward flight of a fixed-wing aircraft. VTOL drones use rotors for vertical ascent and descent, then transition to wing-borne flight for long-range cruise missions. This hybrid approach provides the operational flexibility to take off and land in confined spaces while covering large areas efficiently.}{catD}
\dictentry{VTX}{Video Transmitter, the onboard unit that broadcasts the live camera feed from the drone to the ground-based video receiver. VTX units are characterized by their output power, frequency range, and supported channels, with options ranging from 25 mW for racing to 800 mW or more for long-range flying. Proper VTX power selection is important for both range performance and regulatory compliance.}{catD}
\dictentry{Waypoint Navigation}{An autonomous flight mode where the drone follows a pre-programmed sequence of GPS coordinates, flying from one waypoint to the next without pilot input. Each waypoint specifies a latitude, longitude, altitude, and sometimes additional actions like taking photos or circling. Waypoint navigation enables systematic coverage of large areas and repeatable flight paths essential for mapping and surveying missions.}{catD}
\dictentry{Wildlife Monitoring}{The use of drones to observe, count, and track wildlife populations in their natural habitats without disturbing their behavior. Drones equipped with thermal or high-resolution cameras can survey large areas and detect animals hidden in vegetation or active during nighttime. This non-invasive approach provides more accurate population counts and behavioral data compared to traditional ground surveys.}{catD}
\dictentry{Winch}{A motorized hoisting mechanism mounted on a drone for lowering or raising payloads without the drone needing to land. Winches allow drones to deliver supplies to inaccessible locations, deploy sensors into water, or lower equipment to ground crews. The winch must manage payload weight, cable tension, and descent speed while maintaining the drone's flight stability during the operation.}{catD}
\dictentry{Wind Resistance}{The maximum sustained wind speed in which a drone can maintain stable, controlled flight without excessive power consumption or attitude error. Wind resistance depends on the drone's weight, motor power, aerodynamic profile, and the sophistication of its stabilization algorithms. Operating beyond a drone's wind resistance rating can cause loss of control, battery depletion, and potential flyaway conditions.}{catD}
\dictentry{Yagi Antenna}{A high-gain directional antenna consisting of a driven element, reflector, and one or more directors arranged along a boom. Yagi antennas focus radio energy in a narrow beam, providing significantly more range in the aimed direction than omnidirectional antennas. They are used for long-range FPV video reception and telemetry links where maximum gain in a known direction outweighs the need for broad coverage.}{catD}
\dictentry{Yaw Propeller}{Counter-rotating propellers mounted coaxially that spin in opposite directions to cancel the rotational torque effect on the drone. By having pairs of propellers rotate in opposing directions, the net torque forces balance out, preventing the drone from spinning around its vertical axis. This torque cancellation is fundamental to multirotor stability and eliminates the need for a tail rotor as used on helicopters.}{catD}

\dictentry{Yaw Rate}{The rate of rotation around the drone's vertical axis, measured in degrees per second. Yaw rate is controlled by the pilot's yaw stick input and is stabilized by the flight controller using gyroscope feedback. High yaw rates are used for rapid heading changes in racing and freestyle flying.}{catD}

\dictentry{Zoom Camera}{A camera payload with optical zoom capability, allowing the drone to capture detailed images of distant objects without physically approaching them. Optical zoom preserves image quality unlike digital zoom, making it essential for inspection, surveillance, and search and rescue missions where close approach is impractical or unsafe.}{catD}


